\section*{الفصل \ref{ch:b1:hotspots} --- البقع الساخنة والبراكين داخل الصفائح}

\begin{solution}{exo:b1:hotspots:1}
(أ) $4205 + 5274 = \qty{9479}{m}$، أي نحو \qty{9.5}{km}. فالبركان أثقل من الماء الذي يزيحه: تنحني الصفيحة نحو الأسفل تحت حمله، فتقع قاعدته أعمق من القاع البعيد، وارتفاعه الحقيقي أكبر. (ب) $74000/0.1 = \num{740000}$ سنة بإمداد كيلاويا؛ $74000/0.025 = \num{3.0e6}$ سنة بالمعدّل المتوسط للسلسلة. (ج) $0.12\times19 = \qty{2.3}{km^3}$، نحو $\num{3e-5}$ من ماونا لوا: ثلاثة أجزاء من الألف من واحد في المئة.
\end{solution}

\begin{solution}{exo:b1:hotspots:2}
(أ) $0.92\times\pi/180 = \qty{0.0161}{rad.Myr^{-1}}$. (ب) بالراديان في كل مليون سنة مضروبةً بالكيلومترات، فإن $\omega R\sin\Delta$ تُقاس بالكيلومترات في كل مليون سنة، أي بالمليمترات في السنة: $0.016057\times6371 = 102.3$، إذن $v = 102.3\sin\Delta$، وهذا يعطي \qty{51.1}{mm.yr^{-1}} عند \ang{30}، و\qty{88.6}{mm.yr^{-1}} عند \ang{60}، و\qty{102.3}{mm.yr^{-1}} عند \ang{90}، و\qty{97.3}{mm.yr^{-1}} عند \ang{108.0}. (ج) السرعة متناسبة مع البعد $R\sin\Delta$ عن المحور، وهي أكبر ما تكون على \omterm{def:b1:spherical-astronomy:spherical-triangle}{الدائرة العظمى} الواقعة على \ang{90} من القطب؛ وبعدها تقترب النقطة من المحور ثانية، على جانب القطب المقابل.
\end{solution}

\begin{solution}{exo:b1:hotspots:3}
(أ) الدكن: $\num{e6}/\num{5e5} = \qty{2}{km}$. الأطلسي الأوسط: $\num{2e6}/\num{7e6} = \qty{0.29}{km}$، أي نحو \qty{290}{m}. (ب) المتوسط للدكن قريب من مقطعه البالغ \qty{2.5}{km}: فالمقاطعة هضبة سميكة متصلة. (ج) مساحة الأطلسي الأوسط هي مساحة المنطقة كلها، على أربع قارات، التي توجد عليها بقاياها: حمم وسدود وعتبات الآن مفصولة بمحيط وكثيرة التعرية. والصخر مركّز في بعض المواضع، وغائب عن معظم المساحة، فالمتوسط لا يقول شيئًا عن كومة واحدة.
\end{solution}

\begin{solution}{exo:b1:hotspots:4}
(أ) $\rho_m\alpha\,\Delta T = 3300\times\num{3.2e-5}\times250 =
\qty{26.4}{kg.m^{-3}}$. (ب) يحمل عمود الانتفاخ الكتلة الزائدة $(\rho_m - \rho_w)h = 2300\times690 = \qty{1.587e6}{kg.m^{-2}}$، التي يجب أن توازنها الطبقة الخفيفة: $H = \num{1.587e6}/26.4 = \qty{6.0e4}{m}$، أي نحو \qty{60}{km} من الصخر الحار.
\end{solution}

\begin{solution}{exo:b1:hotspots:5}
(أ) عند A ($\lambda = 0$، $\varphi = 0$)، $\vec e = (0, 1, 0)$ و
$\vec n = (0, 0, 1)$؛ ومع $\sin\ang{270} = -1$ و$\cos\ang{270} = 0$،
$\vec t_A = -\vec e = (0, -1, 0)$. وعند B ($\lambda = \ang{90}$)،
$\vec e = (-1, 0, 0)$، $\vec n = (0, 0, 1)$ و$\vec t_B = (1, 0, 0)$. (ب) $\vec t_A\times\vec t_B = (0, 0, 1)$. (ج) مع $\vec p' = (0, 0, 1)$ و
$\vec r_A = (1, 0, 0)$، $\vec p'\times\vec r_A = (0, 1, 0)$، و
$(\vec p'\times\vec r_A)\cdot\vec t_A = -1 < 0$: القطب هو
$-\vec p' = (0, 0, -1)$، القطب الجنوبي. فالحركة نحو الغرب على امتداد خط الاستواء مع عقارب الساعة منظورةً من فوق القطب الشمالي، فهي عكس عقارب الساعة منظورةً من فوق القطب الجنوبي. (د) عند \ang{90} من القطب، $\omega = 50/6371 =
\qty{7.85e-3}{rad.Myr^{-1}} = \qty{0.45}{\degree\per\Myr}$.
\end{solution}

\begin{solution}{exo:b1:hotspots:6}
(أ) $\sum x_i v_i = 90000 + 240000 + 342000 = \num{672000}$ و
$\sum x_i^2 = \num{7.0e7}$، إذن $\omega = \qty{0.0096}{rad.Myr^{-1}}$، أي
\qty{0.55}{\degree\per\Myr}. (ب) التوقعات هي 28.8 و\num{48.0} و\qty{57.6}{mm.yr^{-1}}؛ والبواقي $+1.2$، $0.0$، و$-0.6$
\unit{mm.yr^{-1}}. (ج) للدوران الصلب السرعة متناسبة مع البعد عن المحور، ونقطة على المحور لا تتحرك: فلا يوجد قطع يُركَّب.
\end{solution}

\begin{solution}{exo:b1:hotspots:7}
(أ) $v = \qty{50e-3}{m}/\qty{3.156e7}{s} = \qty{1.58e-9}{m.s^{-1}}$، إذن
$B = 2300\times\num{3e8}\times\num{1.58e-9} = \qty{1090}{kg.s^{-1}}$ و
$Q = (910/\num{3.2e-5})\times1090 = \qty{3.1e10}{W}$. (ب) ثلث تدفّق هاواي. (ج) يتناسب $B$ مع $A\,v$: فالعمود تحت الصفيحة الأسرع يجب أن يعيد تكوين الانتفاخ نفسه ضعف السرعة، وتدفّقه ضعف ما كان.
\end{solution}

\begin{solution}{exo:b1:hotspots:8}
(أ) $350/(16.1 - 10.3) = 350/5.8 = \qty{60}{km.Myr^{-1}}$، أي
\qty{60}{mm.yr^{-1}}. (ب) يتفق الأثر الفتيّ، \qty{29 \pm 5}{mm.yr^{-1}}، مع \qty{26 \pm 5}{mm.yr^{-1}} للصفيحة ضمن الشكوك؛ أما الجزء القديم فقد كان أسرع من ضعفه، $60 - 26 = 34$، أي نحو سبعة انحرافات معيارية تقريبًا لسرعة الصفيحة. (ج) نعم: الاتجاه شمال \ang{54} شرق هو \omterm{def:b1:spherical-astronomy:horizontal}{السمت} \ang{54}، وهو يعاكس \ang{235} بفارق أقل من درجة. فالبراكين تتحرك فوق الصفيحة في الاتجاه المعاكس لحركة الصفيحة فوق \omterm{def:b1:hotspots:hotspot}{البقعة الساخنة}. (د) الامتداد فرّق الكالديرات القديمة، فزاد المسافات بينها، فزادت السرعة الظاهرية؛ والصفيحة نفسها لم تزدد سرعة.
\end{solution}

\begin{solution}{exo:b1:hotspots:9}
(أ) $\num{9.2e10}/\num{e12} = \qty{0.092}{W.m^{-2}}$، أي \qty{92}{mW.m^{-2}}. (ب) من تسعة إلى ثمانية عشر ضعفًا من الزيادة المرصودة، وقد تكون معدومة. (ج) $L^2/\kappa =
\num{e10}/\num{e-6} = \qty{e16}{s}$، أي نحو \qty{320}{Myr}؛ وتعبر الصفيحة \qty{1000}{km} في $1000/92 = \qty{11}{Myr}$. فللحرارة من الأسفل وقت لتعبر جزءًا صغيرًا فقط من الصفيحة وهي فوق الانتفاخ: فمعظم حرارة العمود تبقى في \omterm{def:b1:layered-earth:layers}{الوشاح} تحت الصفيحة وتُحمل بعيدًا معها، ولا يرى السطح منها إلا القليل.
\end{solution}

\begin{solution}{exo:b1:hotspots:10}
(أ) عشرة أضعاف المصهور، \qty{e7}{km^3}، نصف قطره
$(3V/4\pi)^{1/3} = \qty{134}{km}$. (ب) $\num{e7}/100 = \qty{e5}{km^2}$، قرص
$2\sqrt{\num{e5}/\pi} = \qty{357}{km}$ قطرًا: أي أصغر بثمانية عشر ضعفًا في المساحة من الفخاخ. فالرأس كان أكبر أو ذاب أقل، أو انتشر أوسع وأرقّ، أو تدفّقت الحمم بعيدًا عنه. (ج) \qty{1}{km^3} في السنة، أي عشرة أضعاف الإمداد الحالي لكيلاويا. (د) يجب أن يجمع عمود جديد ما يكفي من صخر حار وخفيف ليشقّ طريقه صعودًا عبر \omterm{def:b1:layered-earth:layers}{الوشاح} اللزج، ويصعد الحجم الكبير أسرع بكثير من الصغير؛ والذيل الذي يتبعه يتدفّق صعودًا في المسار الحار الضعيف الذي فتحه الرأس، ويكفي قناة ضيقة.
\end{solution}

\begin{solution}{exo:b1:hotspots:11}
(أ) $\qty{40}{km.Myr^{-1}}\times(81 - 47)\,\unit{Myr} = \qty{1360}{km}$. (ب) تتحرك الصفيحة بسرعة $(-80\sin\ang{20}, 80\cos\ang{20}) = (-27.4, 75.2)$، \omterm{def:b1:hotspots:hotspot}{والبقعة الساخنة} عند $(0, -40)$: فبالنسبة إلى البقعة تتحرك الصفيحة بمقدار $(-27.4, 115.2)$ \unit{mm.yr^{-1}}، بسرعة قدرها \qty{118.4}{mm.yr^{-1}} نحو $\ang{360} - \arctan(27.4/115.2) = \ang{346.6}$. (ج) يعطي الأثر سرعة أعلى بنسبة \qty{48}{\percent}، وسمتًا أبعد شرقًا بمقدار \ang{6.6}: فالبقعة المنجرفة تسجّل حركة الصفيحة بالنسبة إلى نفسها، لا بالنسبة إلى \omterm{def:b1:layered-earth:layers}{الوشاح} العميق.
\end{solution}

\begin{solution}{exo:b1:hotspots:12}
(أ) $\qty{3e6}{J.m^{-3}.K^{-1}}\times\qty{1300}{K}\times\qty{e15}{m^3} =
\qty{3.9e24}{J}$. (ب) خلال $\qty{3.16e13}{s}$، \qty{1.2e11}{W}: أي قريبًا من الحرارة التي يحملها عمود هاواي اليوم (1.3 مرة)، و\qty{0.26}{\percent} من \qty{47}{TW} حرارة الأرض. (ج) حدّ أدنى: فالرأس يصعد أيضًا بالحرارة الزائدة لتسعة أعشار صخره التي لا تذوب، ويسخّن قاعدة \omterm{def:b1:layered-earth:lithosphere}{الغلاف الصخري}، وقد أُهملت الحرارة الكامنة.
\end{solution}

\begin{solution}{pb:b1:hotspots:1}
\textbf{1.} عند هاواي ($\varphi = \ang{19.0}$، $\lambda = \ang{-155.2}$)، مع $\sin\ang{304} = -0.8290$ و$\cos\ang{304} = 0.5592$:
\[
\begin{aligned}
  \vec r_H &= (-0.8583,\ -0.3966,\ 0.3256), &
  \vec e &= (0.4195,\ -0.9078,\ 0),\\
  \vec n &= (0.2955,\ 0.1366,\ 0.9455), &
  \vec t_H &= (-0.1825,\ 0.8289,\ 0.5287).
\end{aligned}
\]

\textbf{2.} عند لويزفيل ($\varphi = \ang{-53.6}$، $\lambda = \ang{-140.6}$)، مع $\sin\ang{316} = -0.6947$ و$\cos\ang{316} = 0.7193$:
\[
\begin{aligned}
  \vec r_L &= (-0.4586,\ -0.3767,\ -0.8049), &
  \vec e &= (0.6347,\ -0.7727,\ 0),\\
  \vec n &= (-0.6220,\ -0.5109,\ 0.5934), &
  \vec t_L &= (-0.8883,\ 0.1693,\ 0.4269).
\end{aligned}
\]

\textbf{3.} $\lVert\vec t_H\rVert^2 = 0.0333 + 0.6871 + 0.2795 = 0.9999$ و
$\vec t_H\cdot\vec r_H = 0.1566 - 0.3287 + 0.1721 = 0.0000$، في حدود التقريب.

\textbf{4.} $\vec t_H\times\vec t_L = (0.2643, -0.3917, 0.7054)$، وطوله
$0.8491$. والزاوية بين الأثرين هي $\arcsin0.8491 = \ang{58.1}$ (جداؤهما النقطي، $0.528$، موجب، فالزاوية حادّة).

\textbf{5.} لدينا $\vec p' = (0.3113, -0.4613, 0.8308)$. ثم
$(\vec p'\times\vec r_H)\cdot\vec t_H = -0.98 < 0$: القطب هو
$\vec p = (-0.3113, 0.4613, -0.8308)$.

\textbf{6.} $\arcsin(-0.8308) = \ang{-56.18}$ و
$\operatorname{atan2}(0.4613, -0.3113) = \ang{124.01}$: يقع القطب عند \ang{56.2}~ج، \ang{124.0}~ق، في المحيط الجنوبي جنوب أستراليا.

\textbf{7.} $\vec p\cdot\vec r_H = 0.2672 - 0.1830 - 0.2705 = -0.1863$، إذن
$\Delta_H = \ang{100.7}$؛ $\vec p\cdot\vec r_L = 0.1428 - 0.1738 + 0.6687 =
0.6377$، إذن $\Delta_L = \ang{50.4}$.

\textbf{8.} المتجهة $\vec p\times\vec r_L$ عمودية على $\vec p$ وعلى $\vec r_L$، كما أن $\vec t_L$ كذلك (بحسب بناء $\vec p$، ولأن الأثر أفقي)؛ وكلاهما على الخط نفسه، وطولها $\sin\Delta_L$. إذن $(\vec p\times\vec r_L)\cdot\vec t_L = \pm\sin\Delta_L$، والحساب يعطي $+0.770 = \sin\ang{50.4}$: فالصفيحة تتحرك عند لويزفيل على امتداد $+\vec t_L$، في اتجاه امتداد الأثر. وكانت قاعدة الإشارة المطبّقة عند هاواي يمكن أن تعطي الاتجاه الخاطئ عند لويزفيل؛ ولم تفعل، فالاتجاهان يُعادان بدقة تامة، \ang{304} و\ang{316}.

\textbf{9.} $\omega = 67/(6371\sin\ang{50.38}) = 67/4908 =
\qty{0.01365}{rad.Myr^{-1}}$، أي \qty{0.782}{\degree\per\Myr}.

\textbf{10.} $360/0.782 = \qty{460}{Myr}$.

\textbf{11.} مبرهنة جيب التمام بين \ang{56.18}~ج، \ang{124.01}~ق و
\ang{63.78}~ج، \ang{112.48}~ق تعطي \ang{9.5}. ومعدّل الأثرين، \qty{0.78}{\degree\per\Myr}، أقل بنسبة \qty{15}{\percent} من \qty{0.92}{\degree\per\Myr} المركّب على ثلاث عشرة سرعة.

\textbf{12.} سوسايتي \ang{73.6}، بيتكيرن \ang{77.8}، كوب \ang{134.7}.

\textbf{13.} $\omega R = 0.01365\times6371 = \qty{86.98}{mm.yr^{-1}}$، مضروبًا في
$\sin\Delta$: سوسايتي \qty{83.5}{mm.yr^{-1}}، بيتكيرن
\qty{85.0}{mm.yr^{-1}}، كوب \qty{61.8}{mm.yr^{-1}}.

\textbf{14.} يتفق بيتكيرن: \qty{5}{mm.yr^{-1}} أقل من المقاس، في حدود شكّه \qty{15}{mm.yr^{-1}}. وسوسايتي أبطأ بمقدار \qty{25.5}{mm.yr^{-1}}، أي 2.6 من الشكوك، وكوب أسرع بمقدار \qty{18.8}{mm.yr^{-1}}، أي أكثر من ستة: فلا يتفق أيّ منهما.

\textbf{15.} القطب عمودي على متجهتي الأثر كليهما. فإذا كانتا متوازيتين تقريبًا، كان جداؤهما قصيرًا، وخطأ بضع درجات في أيّ اتجاه يديره بزاوية كبيرة: فالقطب محدّد تحديدًا رديئًا.

\textbf{16.} $v = \qty{92e-3}{m}/\qty{3.156e7}{s} = \qty{2.92e-9}{m.s^{-1}}$
و$B = 2300\times\num{4.83e8}\times\num{2.92e-9} = \qty{3240}{kg.s^{-1}}$.

\textbf{17.} $Q = (910/\num{3.2e-5})\times3240 = \qty{9.2e10}{W}$،
\qty{0.20}{\percent} من \qty{47}{TW}.

\textbf{18.} تصنع الحيود $21/67000 = \qty{3.1e-4}{km^3}$ في السنة لكل كيلومتر؛ و$0.1/\num{3.1e-4} = \qty{320}{km}$ من الحيد.

\textbf{19.} $74000/0.1 = \num{740000}$ سنة.

\textbf{20.} $\num{2e6}/\num{81e6} = \qty{0.025}{km^3}$ في السنة: فإمداد كيلاويا وحدها اليوم أربعة أضعاف الناتج المتوسط \omterm{def:b1:hotspots:hotspot}{للبقعة الساخنة} كلها، وهو قريب من أقوى أحواله.

\textbf{21.} $v_H = 86.98\sin\ang{100.73} = 86.98\times0.9825 =
\qty{85.5}{mm.yr^{-1}}$.

\textbf{22.} $92 - 85.5 = \qty{6.5}{mm.yr^{-1}}$، أي 2.2 انحراف معياري من السرعة المقاسة: فالبقعتان لا تتفقان تمامًا على حركة صفيحة صلبة واحدة فوق إطار ثابت.

\textbf{23.} تعطي الأعمار \qty{9.1}{cm} و\qty{8.6}{cm} في السنة، متوسطةً على السلسلة كلها؛ و\qty{9.2}{cm} في السنة عند مورغان ومورغان قريبة من الأولى، وتنبؤ لويزفيل، \qty{8.5}{cm} في السنة، قريب من الثانية. وتقع الأرقام الأربعة كلها ضمن \qty{0.7}{cm} في السنة بعضها من بعض: فالآثار والأعمار والاتجاهات تروي القصة نفسها ضمن نحو \qty{8}{\percent}.

\textbf{24.} بمنزلة عشرية واحدة، يعطي أثر لويزفيل صفيحة المحيط الهادئ فوق \omterm{def:b1:hotspots:hotspot}{البقعة الساخنة} لهاواي السرعة \textbf{\qty[text-series-to-math=true]{8.5}{cm.yr^{-1}}}.
\end{solution}
